Silahkan disesuaikan dengan kebutuhan dan kasusnya masing-masing. Contoh yang bisa diikuti adalah sebagai berikut: 

\section{Alat dan Bahan}
    
\section{Pembagian Data Latih dan Data Uji}
Total data yang digunakan ditunjukkan pada Tabel \ref{tab:total_instruksi}. 

\begin{table}[!h]
    \centering
    \caption{Total Generasi Instruksi}
    \label{tab:total_instruksi}
    \begin{tabular}{lll}
    \toprule
        Deskripsi & Data Paralel & Data Instruksi  \\
    \midrule
        Data Latih  & 53.217 & 212.868 \\
        Data Uji & 13.304 & 13.304 \\
        Experience Replay & - & 1000 \\
    \midrule
        \textbf{Total} & 66.521 & 227.172 \\
    \bottomrule
    \end{tabular}
\end{table}

\section{Pra Pemrosesan}


% \begin{lstlisting}[language=Python, caption=Pelatihan model FastText, label=lst:fasttext_training]
% from gensim.models import FastText
% from nltk.tokenize import word_tokenize
% fasttext_data = []
% fasttext_data.extend(tapaleuk_data['text'].dropna().tolist())
% fasttext_data.extend(jfs_data.dropna().tolist())
% fasttext_data.extend(taxi1500_data.dropna().tolist())
% fasttext_data.extend(puisi_data['text'])
% fasttext_data.extend(pantun_data['text'])
% sentences = [word_tokenize(data.lower()) for data in fasttext_data]
% fasttext_model = FastText(sentences=sentences, vector_size=300, window=5, min_count=2, workers=2)
% \end{lstlisting}


\section{Pelatihan Model}
\subsubsection{Model Training}
 
% \begin{lstlisting}[language=Python, caption=Pelatihan Model, label=lst:pelatihan_model]
% from transformers import Seq2SeqTrainer, Seq2SeqTrainingArguments, DataCollatorForSeq2Seq
% def train():
%     accelerator = Accelerator()
%     model, tokenizer = load_model_and_tokenizer(model_args, logger)
%     model = setup_model_for_training(model, model_args)
%     train_dataset, val_dataset = prepare_datasets(data_args, tokenizer)
%     data_collator = DataCollatorForSeq2Seq(
%         tokenizer, pad_to_multiple_of=8, return_tensors="pt", padding=True)
%     trainer = Seq2SeqTrainer(
%         model=model,
%         args=training_args,
%         train_dataset=train_dataset,
%         eval_dataset=val_dataset,
%         tokenizer=tokenizer,
%         data_collator=data_collator)
%     trainer.train()
%     model = accelerator.unwrap_model(model)
%     model.save_pretrained(
%         training_args.output_dir,
%         is_main_process=accelerator.is_main_process,
%         save_function=accelerator.save)
% \end{lstlisting}

\section{Pengujian Model}


